
\section{Losses and complete results}\label{supp:results}
\subsection{Aggregation and the primary confidence family}
For task $\tau$, arm $m$, and horizon $h\in\{1,4,12\}$, let $v_{\tau h}$ be the retained-A target population variance, floored at $10^{-12}$. The task SMSE is
\[
 L_{\tau m}=\frac13\sum_h\frac{n_D^{-1}\sum_{i\in D}(y_{ih}-\widehat y_{m,ih})^2}{v_{\tau h}},\qquad n_D=1664.
\]
Within repetition $r$, nonlinear loss $L^{\mathrm{NL}}_{r,m}$ averages the six nonlinear law/noise cells equally; linear loss $L^{\mathrm L}_{r,m}$ averages the two linear cells. The independent units are the 32 complete repetition aggregates. Noise levels, horizons, local models, and origins do not increase that sample size.

The four margin-adjusted repetition differences are
\[
 \begin{array}{ll}
 L^{\mathrm{NL}}_{r,P}-0.97L^{\mathrm{NL}}_{r,G},&
 L^{\mathrm{NL}}_{r,K}-0.97L^{\mathrm{NL}}_{r,P},\\
 L^{\mathrm{NL}}_{r,K}-0.97L^{\mathrm{NL}}_{r,R},&
 L^{\mathrm L}_{r,K}-1.05L^{\mathrm L}_{r,P}.
 \end{array}
\]
For each difference sequence $d_r$, define $\bar d=32^{-1}\sum_r d_r$ and $s_d^2=31^{-1}\sum_r(d_r-\bar d)^2$. The upper endpoint is $U=\bar d+t_{31,1-0.05/4}s_d/\sqrt{32}$, where $t_{\nu,p}$ is a Student quantile and $t_{31,0.9875}=2.3555682822$. A criterion is established only for $U<0$. This is one Bonferroni four-comparison family using a paired-$t$ approximation. It is not a distribution-free guarantee. The 3\% and 5\% margins are research conventions rather than deployment-specific costs.

Table~\ref{supp:primary} reports all four primary outcomes. Table~\ref{supp:aggregate} gives their aggregate losses and those of the secondary independent-cap procedure. The K/P gain does not establish that K is preferable to the matched RBF construction, and the linear noninferiority result allows an observed increase in loss. Relative gains are ratios of mean losses, $100(1-\overline L_a/\overline L_b)$, rather than averages of within-repetition ratios.

\begin{table}[htbp]\centering\small
\caption{Complete primary family. The margin-adjusted mean and upper endpoint are in absolute SMSE units. ``Sup.'' denotes superiority and ``NI'' noninferiority. The inferential unit is a repetition aggregate.}\label{supp:primary}
% Generated by build_soft_state_assets.py; do not edit.
% Shared caps; d=L_candidate-factor*L_reference. Four Bonferroni one-sided paired-t upper bounds; met iff upper<0.
\begin{tabular}{@{}lllrrrc@{}}
\toprule
Population & Contrast & Criterion & Gain (\%) & $\bar d$ & Upper bound & Met \\
\midrule
Nonlinear & P vs G & 3\% sup. & 69.20 & \ensuremath{-0.06545} & \ensuremath{-0.06253} & Yes \\
Nonlinear & K vs P & 3\% sup. & 10.46 & \ensuremath{-0.002271} & \ensuremath{-0.001987} & Yes \\
Nonlinear & K vs R & 3\% sup. & 0.03 & 0.0008103 & 0.000981 & No \\
Linear & K vs P & 5\% NI & \ensuremath{-1.11} & \ensuremath{-0.03542} & \ensuremath{-0.03134} & Yes \\
\bottomrule
\end{tabular}
\end{table}
\begin{table}[htbp]\centering\small
\caption{Mean terminal SMSE and repetition standard deviation (SD). Independent cap selection is secondary. SD is not a confidence interval.}\label{supp:aggregate}
% Generated by build_soft_state_assets.py; do not edit.
% D SMSE; SD across 32 complete repetitions, not a confidence interval. Independent caps are secondary.
\begin{tabular}{@{}llrrrr@{}}
\toprule
Cap rule & Arm & Nonlinear mean & SD & Linear mean & SD \\
\midrule
Shared & G & 0.0988703 & 0.00788637 & 0.939042 & 0.0484704 \\
Shared & P & 0.0304507 & 0.00150099 & 0.910921 & 0.0513907 \\
Shared & K & 0.027266 & 0.00123317 & 0.921045 & 0.0526607 \\
Shared & R & 0.027274 & 0.00127755 & 0.915866 & 0.0512334 \\
\addlinespace
Independent & G & 0.0988703 & 0.00788637 & 0.939042 & 0.0484704 \\
Independent & P & 0.0304507 & 0.00150099 & 0.910921 & 0.0513907 \\
Independent & K & 0.0267546 & 0.00120524 & 0.914936 & 0.0511589 \\
Independent & R & 0.0267507 & 0.00129299 & 0.91589 & 0.0521994 \\
\bottomrule
\end{tabular}
\end{table}
\FloatBarrier

\subsection{Law, noise, horizon, and secondary comparisons}
Table~\ref{supp:conditions} retains all eight law/noise conditions under both cap rules. K worsens P on nonlinear autoregression at both noise levels. It improves P on both oscillator laws, but its K/R ordering changes between conditions. Table~\ref{supp:horizons} shows that the nonlinear aggregate K/P benefit is concentrated at horizon twelve. These are descriptive decompositions; they are not additional adjusted subpopulation claims.

The 24 paired descriptive comparisons in Table~\ref{supp:secondary} use ordinary unadjusted 95\% paired-$t$ intervals. For these rows, the difference is candidate loss minus reference loss without a practical-margin factor. Cross-view rows compare the independently selected and shared-cap procedures for the same arm. Neither these intervals nor an independent-cap result replaces a failed primary criterion. In particular, an interval containing zero does not establish equivalence.

\begin{table}[htbp]\centering\small
\caption{Complete descriptive law/noise means over 32 repetitions and three horizons. AR denotes autoregression. The parameter $\eta$ is the added-noise multiplier; oscillator sensor noise is present even at $\eta=0$.}\label{supp:conditions}
% Generated by build_soft_state_assets.py; do not edit.
% Descriptive D SMSE by law/noise; each mean uses all 32 repetitions and averages the three horizons.
\begin{tabular}{@{}llrrrrr@{}}
\toprule
Cap rule & Law & $\eta$ & G & P & K & R \\
\midrule
Shared & Nonlinear AR & 0 & 0.0244281 & 0.0214449 & 0.0226148 & 0.0217499 \\
Shared & Nonlinear AR & 0.1 & 0.0337124 & 0.0329659 & 0.0351019 & 0.0336442 \\
Shared & Van der Pol & 0 & 0.148037 & 0.0151678 & 0.00328116 & 0.00467759 \\
Shared & Van der Pol & 0.1 & 0.149308 & 0.0511417 & 0.0471983 & 0.0477871 \\
Shared & Duffing & 0 & 0.118935 & 0.0158738 & 0.0133303 & 0.0138744 \\
Shared & Duffing & 0.1 & 0.118801 & 0.0461102 & 0.0420697 & 0.0419106 \\
Shared & Linear AR & 0 & 0.937648 & 0.909514 & 0.918951 & 0.914566 \\
Shared & Linear AR & 0.1 & 0.940436 & 0.912328 & 0.923139 & 0.917166 \\
\addlinespace
Independent & Nonlinear AR & 0 & 0.0244281 & 0.0214449 & 0.0225847 & 0.0217757 \\
Independent & Nonlinear AR & 0.1 & 0.0337124 & 0.0329659 & 0.0348346 & 0.0335888 \\
Independent & Van der Pol & 0 & 0.148037 & 0.0151678 & 0.00252973 & 0.00347309 \\
Independent & Van der Pol & 0.1 & 0.149308 & 0.0511417 & 0.046256 & 0.0466229 \\
Independent & Duffing & 0 & 0.118935 & 0.0158738 & 0.0125112 & 0.0133188 \\
Independent & Duffing & 0.1 & 0.118801 & 0.0461102 & 0.0418112 & 0.0417249 \\
Independent & Linear AR & 0 & 0.937648 & 0.909514 & 0.912164 & 0.913393 \\
Independent & Linear AR & 0.1 & 0.940436 & 0.912328 & 0.917709 & 0.918387 \\
\bottomrule
\end{tabular}
\end{table}
\begin{table}[htbp]\centering\small
\caption{Horizon-specific descriptive SMSE. NL and L denote nonlinear and linear panels. Each horizon has a separate fitted direct readout and its own A target-variance normalization.}\label{supp:horizons}
% Generated by build_soft_state_assets.py; do not edit.
% Descriptive horizon-specific D SMSE; NL=nonlinear, L=linear. Per-horizon A target variances normalize loss.
\begin{tabular}{@{}llrrrrrr@{}}
\toprule
Cap rule & Arm & NL $h=1$ & NL $h=4$ & NL $h=12$ & L $h=1$ & L $h=4$ & L $h=12$ \\
\midrule
Shared & G & 0.0300446 & 0.0609004 & 0.205666 & 0.811376 & 0.999827 & 1.00592 \\
Shared & P & 0.0122551 & 0.0192541 & 0.059843 & 0.708544 & 1.00821 & 1.01601 \\
Shared & K & 0.0126386 & 0.0192272 & 0.0499322 & 0.718495 & 1.01811 & 1.02653 \\
Shared & R & 0.0120422 & 0.0186757 & 0.0511041 & 0.710428 & 1.01537 & 1.0218 \\
\addlinespace
Independent & G & 0.0300446 & 0.0609004 & 0.205666 & 0.811376 & 0.999827 & 1.00592 \\
Independent & P & 0.0122551 & 0.0192541 & 0.059843 & 0.708544 & 1.00821 & 1.01601 \\
Independent & K & 0.0123713 & 0.0190399 & 0.0488525 & 0.701798 & 1.01751 & 1.0255 \\
Independent & R & 0.0119662 & 0.018509 & 0.0497769 & 0.706971 & 1.01587 & 1.02483 \\
\bottomrule
\end{tabular}
\end{table}
\begin{table}[htbp]\centering\small
\caption{All planned descriptive paired comparisons. Ordinary 95\% intervals are unadjusted. The difference $\bar\Delta$ is candidate minus reference in absolute SMSE units. Cross-view comparisons use independent minus shared cap selection.}\label{supp:secondary}
% Generated by build_soft_state_assets.py; do not edit.
% Unadjusted descriptive paired intervals; Delta=candidate-reference, with no practical margin. Cross-view rows compare independent minus shared for the same arm.
\begin{tabular}{@{}lllrrrr@{}}
\toprule
Population & Cap rule & Arm contrast & Gain (\%) & $\bar\Delta$ & 95\% lower & 95\% upper \\
\midrule
Nonlinear & Shared & P vs G & 69.20 & \ensuremath{-0.06842} & \ensuremath{-0.07103} & \ensuremath{-0.06581} \\
Nonlinear & Shared & K vs P & 10.46 & \ensuremath{-0.003185} & \ensuremath{-0.00344} & \ensuremath{-0.00293} \\
Nonlinear & Shared & K vs R & 0.03 & \ensuremath{-7.968\times 10^{-6}} & \ensuremath{-0.0001589} & 0.000143 \\
Nonlinear & Shared & R vs P & 10.43 & \ensuremath{-0.003177} & \ensuremath{-0.003391} & \ensuremath{-0.002962} \\
Nonlinear & Independent & P vs G & 69.20 & \ensuremath{-0.06842} & \ensuremath{-0.07103} & \ensuremath{-0.06581} \\
Nonlinear & Independent & K vs P & 12.14 & \ensuremath{-0.003696} & \ensuremath{-0.003961} & \ensuremath{-0.003431} \\
Nonlinear & Independent & K vs R & \ensuremath{-0.01} & \ensuremath{3.864\times 10^{-6}} & \ensuremath{-0.0001309} & 0.0001386 \\
Nonlinear & Independent & R vs P & 12.15 & \ensuremath{-0.0037} & \ensuremath{-0.003929} & \ensuremath{-0.003471} \\
Nonlinear & Ind. vs shared & G & 0.00 & 0 & 0 & 0 \\
Nonlinear & Ind. vs shared & P & 0.00 & 0 & 0 & 0 \\
Nonlinear & Ind. vs shared & K & 1.88 & \ensuremath{-0.0005114} & \ensuremath{-0.0007194} & \ensuremath{-0.0003034} \\
Nonlinear & Ind. vs shared & R & 1.92 & \ensuremath{-0.0005233} & \ensuremath{-0.0007183} & \ensuremath{-0.0003283} \\
\addlinespace
Linear & Shared & P vs G & 2.99 & \ensuremath{-0.02812} & \ensuremath{-0.03159} & \ensuremath{-0.02465} \\
Linear & Shared & K vs P & \ensuremath{-1.11} & 0.01012 & 0.006675 & 0.01357 \\
Linear & Shared & K vs R & \ensuremath{-0.57} & 0.005179 & 0.002502 & 0.007856 \\
Linear & Shared & R vs P & \ensuremath{-0.54} & 0.004945 & 0.003546 & 0.006343 \\
Linear & Independent & P vs G & 2.99 & \ensuremath{-0.02812} & \ensuremath{-0.03159} & \ensuremath{-0.02465} \\
Linear & Independent & K vs P & \ensuremath{-0.44} & 0.004015 & 0.001816 & 0.006214 \\
Linear & Independent & K vs R & 0.10 & \ensuremath{-0.0009539} & \ensuremath{-0.002991} & 0.001083 \\
Linear & Independent & R vs P & \ensuremath{-0.55} & 0.004969 & 0.003123 & 0.006815 \\
Linear & Ind. vs shared & G & 0.00 & 0 & 0 & 0 \\
Linear & Ind. vs shared & P & 0.00 & 0 & 0 & 0 \\
Linear & Ind. vs shared & K & 0.66 & \ensuremath{-0.006109} & \ensuremath{-0.008889} & \ensuremath{-0.003328} \\
Linear & Ind. vs shared & R & \ensuremath{-2.65\times 10^{-3}} & \ensuremath{2.431\times 10^{-5}} & \ensuremath{-0.00188} & 0.001929 \\
\bottomrule
\end{tabular}
\end{table}
\FloatBarrier

\subsection{Selected dimensions and numerical coverage}
The P-selected shared caps $1,2,4,8,12$ occur in $31,104,75,32,14$ tasks. Table~\ref{supp:selection} reports all arms and both cap rules, and Table~\ref{supp:selectionpanels} separates nonlinear and linear panels. G's gates and predictions are identical across caps; its independent selection of cap one is tie-breaking, not evidence of optimal rank-one dynamics. P's shared and independent choices coincide by definition.

Shared-bank widths are 14 for G, 28--182 for P, and 196--350 for K and R. Centered A feature ranks range from 9--11, 23--166, 186--334, and 180--334, respectively. These are feature-matrix ranks, not local PCA ranks or independent sample sizes. Five of 17,920 local-model/cap records use the rank-zero wrapper: one singleton model in one task under all five caps. No positive-rank model requires a projected-range reduction. The retained B scores cover all 76,800 cap/arm/penalty/horizon entries, with no discarded task and no recorded SVD-solver recovery.

\begin{table}[htbp]\centering\small
\caption{Cap-selection counts over 256 tasks and ranges of feature width and centered A feature rank. Width excludes the intercept. Shared rows represent the same P-selected bank.}\label{supp:selection}
% Generated by build_soft_state_assets.py; do not edit.
% B-selected cap counts across all 256 task conditions (not independent inferential units); width excludes intercept; centered feature rank on A.
\begin{tabular}{@{}llrrrrrcc@{}}
\toprule
Cap rule & Arm & 1 & 2 & 4 & 8 & 12 & Width range & Rank range \\
\midrule
Shared & G & 31 & 104 & 75 & 32 & 14 & 14 & 9--11 \\
Shared & P & 31 & 104 & 75 & 32 & 14 & 28--182 & 23--166 \\
Shared & K & 31 & 104 & 75 & 32 & 14 & 196--350 & 186--334 \\
Shared & R & 31 & 104 & 75 & 32 & 14 & 196--350 & 180--334 \\
\addlinespace
Independent & G & 256 & 0 & 0 & 0 & 0 & 14 & 9--11 \\
Independent & P & 31 & 104 & 75 & 32 & 14 & 28--182 & 23--166 \\
Independent & K & 67 & 57 & 73 & 42 & 17 & 196--350 & 181--334 \\
Independent & R & 59 & 75 & 89 & 16 & 17 & 196--350 & 177--334 \\
\bottomrule
\end{tabular}
\end{table}
\begin{table}[htbp]\centering\small
\caption{Cap counts by panel. NL has 192 law/noise/repetition tasks and L has 64; each panel still has 32 independent repetition units.}\label{supp:selectionpanels}
% Generated by build_soft_state_assets.py; do not edit.
% Descriptive cap counts: 192 nonlinear and 64 linear task conditions, each arising from 32 repetitions.
\begin{tabular}{@{}lllrrrrr@{}}
\toprule
Population & Cap rule & Arm & 1 & 2 & 4 & 8 & 12 \\
\midrule
Nonlinear & Shared & G & 25 & 89 & 61 & 14 & 3 \\
Nonlinear & Shared & P & 25 & 89 & 61 & 14 & 3 \\
Nonlinear & Shared & K & 25 & 89 & 61 & 14 & 3 \\
Nonlinear & Shared & R & 25 & 89 & 61 & 14 & 3 \\
Nonlinear & Independent & G & 192 & 0 & 0 & 0 & 0 \\
Nonlinear & Independent & P & 25 & 89 & 61 & 14 & 3 \\
Nonlinear & Independent & K & 66 & 55 & 50 & 18 & 3 \\
Nonlinear & Independent & R & 56 & 64 & 66 & 4 & 2 \\
\addlinespace
Linear & Shared & G & 6 & 15 & 14 & 18 & 11 \\
Linear & Shared & P & 6 & 15 & 14 & 18 & 11 \\
Linear & Shared & K & 6 & 15 & 14 & 18 & 11 \\
Linear & Shared & R & 6 & 15 & 14 & 18 & 11 \\
Linear & Independent & G & 64 & 0 & 0 & 0 & 0 \\
Linear & Independent & P & 6 & 15 & 14 & 18 & 11 \\
Linear & Independent & K & 1 & 2 & 23 & 24 & 14 \\
Linear & Independent & R & 3 & 11 & 23 & 12 & 15 \\
\bottomrule
\end{tabular}
\end{table}
\FloatBarrier

\subsection{Resource measurements and computational verification}
Table~\ref{supp:resources} measures the complete candidate search, including rejected candidates and shared geometry work. Fit wall time totals 931.88 seconds, and D prediction wall time totals 671.84 seconds. These are sums over worker tasks, not elapsed time for the three-worker stage. Peak resident set size (RSS) is measured per worker; its maximum is 156.99 MiB, not aggregate concurrent memory. Portable NPZ fitted states occupy 177.93 MiB, including all candidates rather than a minimal deployed model. C diagnostics are included in fitting costs. Preparation, audit, analysis, document production, and compression are outside these measurements.

The retained audit reports 2,563,814 checks across 256 tasks. It independently refits 25,600 B penalty-path arrays using Cholesky calculations, rather than the production SVD calculation, and replays 2,048 selected C and 2,048 selected D forecast arrays. Maximum absolute discrepancies are $1.20\times10^{-11}$ for B, $2.43\times10^{-12}$ for C, and $5.59\times10^{-12}$ for D. The declared replay tolerances are $2\times10^{-7}$ relative and $2\times10^{-9}$ absolute. The audit also checks saved references, local-rank rules, gates, block transformations, selected penalties, and primary statistics. Source and prefix hashes are checked exactly; independently reconstructed numerical arrays use the stated tolerances.

This coverage concerns the retained cohort. The audit does not independently reimplement partition fitting or physical-law integration. It is a computational verification using a distinct calculation for specified components, not a replication by an independent research group. Source publication and stage seals document ordering, but terminal targets are physically present after D preparation. They are not protected by cryptographic isolation.

\begin{table}[htbp]\centering\small
\caption{Whole-search resource measurements. Wall times are worker sums; peak RSS is per worker. State sizes include all candidate models. Shared computations prevent these totals from identifying isolated K versus R deployment costs.}\label{supp:resources}
% Generated by build_soft_state_assets.py; do not edit.
% Complete search/all four arms. Summed worker wall time is not parallel elapsed time; RSS is not concurrent aggregate memory; state sizes retain all candidates.
\begin{tabular}{@{}lrrr@{}}
\toprule
Measured quantity & Task sum & Task median & Task maximum \\
\midrule
Whole fit wall time (s) & 931.88 & 3.63 & 4.75 \\
Whole fit CPU time (s) & 929.02 & 3.62 & 4.65 \\
D prediction wall time (s) & 671.84 & 2.64 & 3.44 \\
D prediction CPU time (s) & 670.60 & 2.64 & 3.43 \\
Peak worker RSS (MiB) & --- & 155.88 & 156.99 \\
All-candidate NPZ states (MiB) & 177.93 & 0.69 & 0.70 \\
All-candidate JSON states (MiB) & 88.24 & 0.34 & 0.34 \\
\bottomrule
\end{tabular}
\end{table}
\FloatBarrier
